% v3-figures-ru/federation.tex: three authorities as a triangle inscribed in a circle.
% The sides are the trust relations: Б attests А, А attests В, and the side between Б and В
% is the refused transitive edge. The holder and the relying party stand inside.
\begin{tikzpicture}[
  auth/.style={layerbox, text width=4.2cm, minimum height=2.35cm, align=center, font=\scriptsize},
]
\begin{scope}[on background layer]
  \draw[gold, line width=0.6pt, dash pattern=on 3pt off 2pt] (0,0) circle (7.00);
\end{scope}
\node[auth] (A) at (90:7.00) {{\bfseries\small Орган А}\\[2pt]своя база данных, свой корень МЛ-ДСА-87; публикует опись, контрольную точку эпохи, ленту отзывов};
\node[auth] (B) at (330:7.00) {{\bfseries\small Орган Б}\\[2pt]своя база данных, свой корень МЛ-ДСА-65; публикует те же три объекта};
\node[auth, fill=steelgray!85, draw=steelgray] (C) at (210:7.00) {{\bfseries\small Орган В}\\[2pt]доверенный для А,\\но не для Б};
\draw[flow, line width=1.4pt] (B) -- (A);
\node[font=\tiny, text=navydark, text width=2.7cm, align=center] at (30:4.95) {\textbf{Б аттестует ключ А} в контексте 3, до такой-то даты\\(\code{АттестацияДоверия})};
\draw[flowfaint] (A) -- (C);
\node[font=\tiny\itshape, text=steelgray, text width=2.5cm, align=center] at (150:4.95) {А аттестует В в каком-то контексте};
\draw[refuse, dash pattern=on 4pt off 3pt] (B) -- node[below=3pt, font=\scriptsize\bfseries\color{classmark}, text width=7.6cm, align=center] {никакого замыкания: доверие Б к А никогда не влечёт доверия Б к В} (C);
\node[humanbox, text width=3.3cm, align=center, font=\scriptsize\bfseries] (h) at (0,1.75) {держатель удостоверения А};
\node[extbox, text width=4.3cm, align=center, font=\scriptsize] (rp) at (0,-1.05) {{\bfseries доверяющая сторона, доверяющая Б}\\[2pt]решает автономно: удостоверение А подлинно, опись Б подлинна, свежа и доверена, она аттестует ключ А в контексте 3, а лента А удостоверения не перечисляет};
\draw[flow] (h) -- node[right=2pt, font=\tiny\itshape, text=steelgray] {предъявляет в контексте 3} (rp);
\draw[flowfaint] (B.north west) to[out=150, in=0] node[pos=0.5, above=1pt, font=\tiny\itshape, text=steelgray] {забирает опись Б} (rp.east);
\node[faint, text width=12.5cm, align=center, anchor=north] at (0,-7.35) {В НИ это выполняется как два независимых экземпляра поверх ХТТП, у каждого своя база данных и настоящий ключ, при смешанной паре наборов параметров. Два экземпляра одной кодовой базы, запущенные одним автором, доказывают протокол; институциональной независимости они не доказывают.};
\end{tikzpicture}
